% 4-5-dm.tex
%  (section opener -- the substance lives in 4-5-1 and 4-5-2)
%  Restates the two declared families from 3-5 BEFORE any p-value is quoted.
%  Rules: decimals in \lr{}; integers Persian; cite only from references.bib.

\section{آزمون دیبولد-ماریانو}
\label{sec:4-5}

جدول‌های \crefrange{tab:hourly-results}{tab:direct-vs-agg} اختلاف‌های عددی را نشان دادند، اما اختلاف
عددی به‌تنهایی چیزی را اثبات نمی‌کند. دو مدل که بر ۷۲۸ مبدأ پیش‌بینی، میانگین خطایی
با فاصله‌ی چند صدم یورو دارند، ممکن است صرفاً بازتاب همان نمونه‌ی خاص باشند. آزمون
دیبولد-ماریانو دقیقاً همین را می‌سنجد: آیا سری اختلافِ زیانِ دو مدل، میانگینی
به‌طور نظام‌مند متفاوت از صفر دارد؟

سه قید که در \cref{sec:3-5} تثبیت شد و در سراسر این بخش برقرار است:

\begin{itemize}
  \item تمام \lr{p}-مقدارها با تصحیح \lr{HAC} برای خودهمبستگی و ناهمسانیِ
        واریانس محاسبه شده‌اند. \lr{p}-مقدارهای تصحیح‌نشده در هیچ جای این فصل
        نقل نمی‌شوند.
  \item آزمون‌ها یک‌طرفه‌اند. در هر جدول، \lr{p} احتمالِ آن است که مدلِ سطر
        \emph{بهتر} از مدلِ ستون نباشد؛ پس مقدار کوچک یعنی سطر بهتر است.
  \item دو خانواده‌ی آزمون از پیش اعلام شده‌اند و اینجا از هم جدا می‌مانند.
\end{itemize}

خانواده‌ی نخست \lr{—} خانواده‌ی تأییدی \lr{—} تنها یک فرضیه دارد: اینکه
وزن‌دهیِ رژیم‌محور بر وزن‌دهیِ ایستا برتری دارد. ثبتِ آن در دفتر تصمیم‌های پروژه
به یازدهم ژوئیه ۲۰۲۶ بازمی‌گردد، یعنی زمانی که هنوز هیچ ترکیبی بر پنجره‌ی
آزمون امتیازدهی نشده بود. یک فرضیه‌ی از پیش تعیین‌شده بدون تصحیح در سطح
\lr{0.05} آزموده می‌شود.

خانواده‌ی دوم \lr{—} خانواده‌ی اکتشافی \lr{—} کل ماتریس ۲۱تاییِ مقایسه‌های
دوبه‌دوی \cref{sec:4-5-1} است. برای این خانواده، هم \lr{p}-مقدار خام و هم
\lr{p}-مقدار تصحیح‌شده با روش هولم\lr{–}بونفرونی گزارش می‌شود، و هر ادعای معناداری
تنها بر پایه‌ی ستون تصحیح‌شده بیان می‌گردد. گزارش‌کردن تنها ستون خام، در ماتریسی با
۲۱ آزمون، به‌طور متوسط یک برنده‌ی کاذب تولید می‌کند.

بخش \cref{sec:4-5-1} مدل‌های این پژوهش را با یکدیگر می‌سنجد و \cref{sec:4-5-2}
آن‌ها را در برابر پیش‌بینی‌های منتشرشده‌ی لاگو و همکاران \cite{lago_forecasting_2021}
قرار می‌دهد.
